\documentclass{article}
\usepackage{geometry}
\geometry{margin=1.5cm, vmargin={0pt,1cm}}
\setlength{\topmargin}{-1cm}
\setlength{\headheight}{12.5pt}
\setlength{\paperheight}{29.7cm}
\setlength{\textheight}{25.3cm}

% useful packages.
\usepackage[UTF8]{ctex}
\usepackage{fancyhdr}
\usepackage{layout}
\usepackage{tikz}

% import common macros
% % % % % % % % % % % % % % % % % % % % % % % % % % % % % % % %
% Common Macros for LaTeX
% % % % % % % % % % % % % % % % % % % % % % % % % % % % % % % %

% Useful packages
\usepackage{amsfonts}
\usepackage{amsmath}
\usepackage{amssymb}
\usepackage{amsthm}
\usepackage{enumerate}
\usepackage{graphicx}
\usepackage{multicol}
\usepackage[ruled,linesnumbered]{algorithm2e}

% Math operators and common symbols
\newcommand{\dif}{\mathrm{d}}
\newcommand{\innerProd}[1]{\left\langle #1 \right\rangle}
\newcommand{\difFrac}[2]{\frac{\dif #1}{\dif #2}}
\newcommand{\pdfFrac}[2]{\frac{\partial #1}{\partial #2}}
\newcommand{\T}{\mathrm{T}}
\newcommand{\norm}[1]{\left\| #1 \right\|}
\newcommand{\abs}[1]{\left| #1 \right|}

% Vector and Matrix shortcuts
\newcommand{\vecTwo}[2]{
  \begin{bmatrix}
    #1 \\
    #2
\end{bmatrix}}
\newcommand{\vecThree}[3]{
  \begin{bmatrix}
    #1 \\
    #2 \\
    #3
\end{bmatrix}}
\newcommand{\vecTwoH}[2]{
  \begin{bmatrix}
    #1 & #2
  \end{bmatrix}
}
\newcommand{\vecThreeH}[3]{
  \begin{bmatrix}
    #1 & #2 & #3
  \end{bmatrix}
}

% Common fields
\newcommand{\R}{\mathbb{R}}
\newcommand{\C}{\mathbb{C}}
\newcommand{\Z}{\mathbb{Z}}
\newcommand{\N}{\mathbb{N}}

% Solutions and proofs
\newenvironment{solution}{
\begin{proof}[解]}{
\end{proof}}

% Utility to get environment variables (requires lualatex)
\newcommand{\getEnv}[2]{\directlua{tex.print(os.getenv("#1") or "#2")}}


% Name and uID
\newcommand{\name}{\getEnv{NAME}{NAME}}
\newcommand{\uid}{\getEnv{UID}{UID}}
\newcommand{\course}{Real Analysis}
\newcommand{\hwNum}{15}

\begin{document}

\pagestyle{fancy}
\fancyhead{}
\lhead{\zhstr{\name} (\uid)}
\chead{\course\ \#\hwNum}
\rhead{\today}

% ==============================================================================
% Page 119, Exercise 29
% ==============================================================================
\section{P119 Exercise 29}

设 $\alpha>0,\ \beta>0$,
\[
  f(x)=x^\alpha\sin x^{-\beta},\quad 0<x\le 1,\qquad f(0)=0.
\]
求证: 当且仅当 $\alpha>\beta$ 时 $f(x)$ 绝对连续.

\begin{solution}
  先设 $\alpha>\beta$. 对 $x>0$,
  \[
    f'(x)=\alpha x^{\alpha-1}\sin x^{-\beta}
      -\beta x^{\alpha-\beta-1}\cos x^{-\beta}.
  \]
  因为 $\alpha>0$ 且 $\alpha-\beta>0$,
  \[
    |f'(x)|\le \alpha x^{\alpha-1}
      +\beta x^{\alpha-\beta-1}\in L([0,1]).
  \]
  对任意 $0<\varepsilon<x\le1$, Newton--Leibniz 公式给出
  \[
    f(x)-f(\varepsilon)=\int_\varepsilon^x f'(t)\,\dif t.
  \]
  令 $\varepsilon\to0^+$, 由 $f(\varepsilon)\to0=f(0)$ 及
  $f'\in L([0,1])$ 可得
  \[
    f(x)=\int_0^x f'(t)\,\dif t.
  \]
  故 $f$ 是 $L^1$ 函数的不定积分, 在 $[0,1]$ 上绝对连续.

  反之, 设 $\alpha\le\beta$. 令
  \[
    x_n=\left(\frac{\pi}{2}+n\pi\right)^{-1/\beta},\qquad n=0,1,2,\ldots.
  \]
  则 $x_n\downarrow0$ 且
  \[
    f(x_n)=(-1)^n x_n^\alpha.
  \]
  对由 $0,x_N,x_{N-1},\ldots,x_0$ 构成的分割,
  \[
    V_0^1(f)\ge \sum_{n=0}^{N-1}|f(x_{n+1})-f(x_n)|
    =\sum_{n=0}^{N-1}(x_{n+1}^\alpha+x_n^\alpha).
  \]
  而 $x_n^\alpha\asymp n^{-\alpha/\beta}$, 当
  $\alpha/\beta\le1$ 时该级数发散, 故 $V_0^1(f)=\infty$.
  绝对连续函数必有界变差, 所以此时 $f$ 不绝对连续.
  故 $f$ 绝对连续当且仅当 $\alpha>\beta$.
\end{solution}

% ==============================================================================
% Page 119, Exercise 31
% ==============================================================================
\section{P119 Exercise 31}

设 $f\in L([a,b])$. 若
\[
  \int_a^b|f(x+h)-f(x)|\,\dif x=o(|h|),
\]
求证: 除一零测集外 $f$ 为常数.

\begin{solution}
  令
  \[
    F(x)=\int_a^x f(t)\,\dif t,\qquad
    E=\{x\in(a,b):F'(x)=f(x)\}.
  \]
  由 Lebesgue 微分定理, $m([a,b]\setminus E)=0$.

  任取 $x,y\in E,\ x<y$. 当 $|h|$ 充分小时,
  \begin{align*}
    &\frac{F(y+h)-F(y)}{h}
      -\frac{F(x+h)-F(x)}{h}\\
    &\quad=\frac1h\left(\int_y^{y+h}f(t)\,\dif t
      -\int_x^{x+h}f(t)\,\dif t\right)\\
    &\quad=\frac1h\int_x^y\bigl(f(t+h)-f(t)\bigr)\,\dif t.
  \end{align*}
  从而
  \[
    \left|
      \frac{F(y+h)-F(y)}{h}
      -\frac{F(x+h)-F(x)}{h}
    \right|
    \le \frac1{|h|}\int_a^b|f(t+h)-f(t)|\,\dif t\longrightarrow0.
  \]
  令 $h\to0$, 由 $x,y\in E$ 得 $f(y)-f(x)=0$.
  故 $f$ 在 $E$ 上取常值, 即几乎处处为常数.
\end{solution}

% ==============================================================================
% Page 119, Exercise 33
% ==============================================================================
\section{P119 Exercise 33}

设对每一 $0<x<1$, $f$ 在 $[0,x]$ 上绝对连续, 在 $[x,1]$ 上有界变差,
而且 $f$ 在 $x=1$ 连续. 求证: $f$ 在 $[0,1]$ 上绝对连续.

\begin{solution}
  固定 $c\in(0,1)$. 因为 $f\in BV([c,1])$, 其变差函数
  \[
    V(t)=V_c^t(f),\qquad c\le t\le1
  \]
  单调递增. 对任意 $t<1$, $f$ 在 $[0,t]$ 上绝对连续, 故也在
  $[c,t]$ 上绝对连续, 从而
  \[
    V(t)=\int_c^t|f'(s)|\,\dif s.
  \]
  令 $L=\lim_{t\to1^-}V(t)$. 由 $f$ 在 $1$ 连续可知, 在分割中把端点
  $1$ 换成充分接近 $1$ 的点不会改变变差的上确界, 故
  \[
    V_c^1(f)=L.
  \]
  所以
  \[
    V_t^1(f)=V_c^1(f)-V_c^t(f)\longrightarrow0
    \qquad(t\to1^-).
  \]

  任给 $\varepsilon>0$, 取 $t\in(c,1)$ 使
  $V_t^1(f)<\varepsilon/2$. 又因 $f$ 在 $[0,t]$ 上绝对连续,
  存在 $\delta>0$, 使得 $[0,t]$ 中任意有限个互不相交区间
  $\{(a_k,b_k)\}$ 满足 $\sum_k(b_k-a_k)<\delta$ 时,
  \[
    \sum_k|f(b_k)-f(a_k)|<\frac{\varepsilon}{2}.
  \]
  对 $[0,1]$ 中任意满足同样长度条件的有限个互不相交区间,
  在 $t$ 处分割跨过 $t$ 的区间. 位于 $[0,t]$ 中的各段增量之和小于
  $\varepsilon/2$, 位于 $[t,1]$ 中的各段增量之和不超过
  $V_t^1(f)<\varepsilon/2$. 故总增量之和小于 $\varepsilon$.
  故 $f$ 在 $[0,1]$ 上绝对连续.
\end{solution}

% ==============================================================================
% Page 119, Exercise 35
% ==============================================================================
\section{P119 Exercise 35}

设 $f\in L([0,1])$, $f$ 在 $0$ 的一个邻域中有界.
\begin{enumerate}
  \renewcommand{\labelenumi}{(\roman{enumi})}
  \item 证明: 对任何 $n\ge1$, $f(x^n)\in L([0,1])$;
  \item 若 $f$ 单调, 求极限 $\displaystyle\lim_{n\to\infty}\int_0^1f(x^n)\,\dif x$.
\end{enumerate}

\begin{solution}
  (i)\; 作代换 $t=x^n$, 得
  \[
    \int_0^1|f(x^n)|\,\dif x
    =\frac1n\int_0^1|f(t)|t^{1/n-1}\,\dif t.
  \]
  取 $\delta\in(0,1)$ 及 $M>0$, 使得 $|f(t)|\le M$ 于
  $(0,\delta)$. 则
  \[
    \frac1n\int_0^\delta|f(t)|t^{1/n-1}\,\dif t
    \le M\delta^{1/n}<\infty.
  \]
  在 $[\delta,1]$ 上, $t^{1/n-1}\le\delta^{1/n-1}$, 故
  \[
    \frac1n\int_\delta^1|f(t)|t^{1/n-1}\,\dif t
    \le\frac{\delta^{1/n-1}}n\int_\delta^1|f(t)|\,\dif t<\infty.
  \]
  故 $f(x^n)\in L([0,1])$.

  (ii)\; 若 $f$ 单调, 则 $f$ 在 $[0,1]$ 上有界, 且右极限
  $f(0+)=\lim_{t\to0^+}f(t)$ 存在. 对每个 $0\le x<1$,
  $x^n\to0$, 从而
  \[
    f(x^n)\longrightarrow f(0+).
  \]
  又 $\{f(x^n)\}$ 被一个常数控制, 由控制收敛定理,
  \[
    \lim_{n\to\infty}\int_0^1f(x^n)\,\dif x=f(0+).
    \qed
  \]
\end{solution}

% ==============================================================================
% Page 135, Exercise 1
% ==============================================================================
\section{P135 Exercise 1}

设 $f\in L^p(\mathbf R),\ 1\le p<\infty$. 求证: 当 $\lambda\to\infty$ 时,
\[
  \int_{\{|f|>\lambda\}}|f(x)|^p\,\dif x\to0,\qquad
  \int_{\{|x|>\lambda\}}|f(x)|^p\,\dif x\to0.
\]

\begin{solution}
  因为 $|f|^p\in L^1(\mathbf R)$, 且当 $\lambda\to\infty$ 时
  \[
    \chi_{\{|f|>\lambda\}}(x)\to0,\qquad
    \chi_{\{|x|>\lambda\}}(x)\to0
  \]
  对 a.e.\ $x$ 成立. 由控制收敛定理,
  \[
    \int_{\{|f|>\lambda\}}|f|^p\to0,\qquad
    \int_{\{|x|>\lambda\}}|f|^p\to0. \qed
  \]
\end{solution}

% ==============================================================================
% Page 135, Exercise 2
% ==============================================================================
\section{P135 Exercise 2}

设 $f\in L^p(\mathbf R),\ 1\le p<\infty$. 求证: 当
$\lambda\to0^+$ 及 $\lambda\to+\infty$ 时皆有
\[
  \lambda^p m(\{|f|>\lambda\})\to0.
\]

\begin{solution}
  记 $E_\lambda=\{x:|f(x)|>\lambda\}$.

  当 $\lambda\to+\infty$ 时,
  \[
    0\le\lambda^p m(E_\lambda)
    \le\int_{E_\lambda}|f(x)|^p\,\dif x\longrightarrow0,
  \]
  其中最后一步由上一题得到.

  当 $\lambda\to0^+$ 时, 任给 $\varepsilon>0$, 由 $f\in L^p(\mathbf R)$
  可取 $A>0$ 使
  \[
    \int_{\{|x|>A\}}|f(x)|^p\,\dif x<\varepsilon.
  \]
  将 $E_\lambda$ 分为
  \[
    E_\lambda^{(1)}=E_\lambda\cap\{|x|\le A\},\qquad
    E_\lambda^{(2)}=E_\lambda\cap\{|x|>A\}.
  \]
  则
  \[
    \lambda^p m(E_\lambda^{(1)})\le2A\lambda^p\to0,
  \]
  且
  \[
    \lambda^p m(E_\lambda^{(2)})
    \le\int_{\{|x|>A\}}|f(x)|^p\,\dif x<\varepsilon.
  \]
  因而
  $\limsup_{\lambda\to0^+}\lambda^p m(E_\lambda)\le\varepsilon$.
  令 $\varepsilon\to0^+$ 即得. \qed
\end{solution}

% ==============================================================================
% Page 135, Exercise 5
% ==============================================================================
\section{P135 Exercise 5}

设 $f_k\in L^{p_k}(E),\ k=1,2,\ldots,n$, 其中 $p_k>1$,
\[
  \frac1{p_1}+\cdots+\frac1{p_n}=1.
\]
求证:
\[
  \int_E\prod_{k=1}^n|f_k(x)|\,\dif x
  \le\prod_{k=1}^n\|f_k\|_{p_k}.
\]

\begin{solution}
  对 $n$ 作归纳. $n=2$ 时结论就是 Hölder 不等式.

  假设结论对 $n-1$ 个函数成立. 令
  \[
    \frac1q=\sum_{k=1}^{n-1}\frac1{p_k}
      =1-\frac1{p_n},
  \]
  则 $q$ 与 $p_n$ 共轭. 因为
  \[
    \sum_{k=1}^{n-1}\frac{q}{p_k}=1,
  \]
  对函数 $|f_k|^q$ 及指数 $p_k/q$ 应用归纳假设, 得
  \[
    \int_E\prod_{k=1}^{n-1}|f_k|^q
    \le\prod_{k=1}^{n-1}
      \left(\int_E|f_k|^{p_k}\right)^{q/p_k}.
  \]
  故
  \[
    \left\|\prod_{k=1}^{n-1}|f_k|\right\|_q
    \le\prod_{k=1}^{n-1}\|f_k\|_{p_k}.
  \]
  再对 $\prod_{k=1}^{n-1}|f_k|$ 与 $|f_n|$ 应用通常的 Hölder 不等式,
  \[
    \int_E\prod_{k=1}^n|f_k|
    \le
    \left\|\prod_{k=1}^{n-1}|f_k|\right\|_q\|f_n\|_{p_n}
    \le\prod_{k=1}^n\|f_k\|_{p_k}.
    \qed
  \]
\end{solution}

% ==============================================================================
% Page 135, Exercise 6
% ==============================================================================
\section{P135 Exercise 6}

设 $f$ 和 $g$ 在 $[a,b]$ 上非负可测而且 $g\in L([a,b])$. 求证:
\[
  \left(\int_a^b f(x)g(x)\,\dif x\right)^p
  \le \|g\|_1^{p-1}\int_a^b f^p(x)g(x)\,\dif x,
  \qquad 1\le p<\infty.
\]

\begin{solution}
  当 $p=1$ 时等号成立. 设 $p>1$, 令 $q=p/(p-1)$.
  将被积函数写成
  \[
    f g=(f g^{1/p})g^{1/q}.
  \]
  由 Hölder 不等式,
  \[
    \int_a^b fg
    \le
    \left(\int_a^b f^p g\right)^{1/p}
    \left(\int_a^b g\right)^{1/q}.
  \]
  两边取 $p$ 次方, 并注意到 $p/q=p-1$, 即得
  \[
    \left(\int_a^b fg\right)^p
    \le\left(\int_a^b g\right)^{p-1}\int_a^b f^p g
    =\|g\|_1^{p-1}\int_a^b f^p g.
    \qed
  \]
\end{solution}

\end{document}
